\section{The Rules, the Normal Boxes, and the Size of the Development}\label{sec:appendix}

\begin{center}
\begin{minipage}[t]{0.42\textwidth}\centering\scriptsize
  \begin{tabular}{@{}r@{\;\,}r@{\,}c@{\,}l@{}}
    order-S & \ufig[0.48]{pv0-order-S-l} & $=$ & \ufig[0.48]{pv0-order-S-r}\\[1pt]
    order-H & \ufig[0.48]{pv0-order-H-l} & $=$ & \ufig[0.48]{pv0-order-H-r}\\[1pt]
    M-power & \ufig[0.48]{pv0-M-power-l} & $=$ & \ufig[0.48]{pv0-M-power-r}\\[1pt]
    semi-MR & \ufig[0.48]{pv0-semi-MR-l} & $=$ & \ufig[0.48]{pv0-semi-MR-r}\\[1pt]
    order-CZ & \ufig[0.48]{pv0-order-CZ-l} & $=$ & \ufig[0.48]{pv0-order-CZ-r}\\[1pt]
    order-Ex & \ufig[0.48]{pv0-order-Ex-l} & $=$ & \ufig[0.48]{pv0-order-Ex-r}\\[1pt]
    order-SH & \ufig[0.48]{pv0-order-SH-l} & $=$ & \ufig[0.48]{pv0-order-SH-r}\\
  \end{tabular}
\end{minipage}\hfill
\begin{minipage}[t]{0.55\textwidth}\centering\scriptsize
  \begin{tabular}{@{}r@{\;\,}l@{}}
    comm-CZ-S↑ & \ufig[0.48]{pv0-comm-CZ-Sup}\\[1pt]
    semi-M↑CZ & \ufig[0.48]{pv0-semi-Mup-CZ}\\[1pt]
    semi-Ex-S↑ & \ufig[0.48]{pv0-semi-Ex-Sup}\\[1pt]
    semi-Ex-H↑ & \ufig[0.48]{pv0-semi-Ex-Hup}\\[1pt]
    comm-HHSHHS & \ufig[0.4]{pv0-comm-HHSHHS-l} $=$ \ufig[0.4]{pv0-comm-HHSHHS-r}\\
  \end{tabular}
\end{minipage}\\[3pt]
\scriptsize
\begin{tabular}{@{}r@{\;\,}l@{\qquad}r@{\;\,}l@{}}
  blake-c12 & \ufig[0.48]{pv0-blake-c12} & yang-baxter & \ufig[0.48]{pv0-yang-baxter}\\[2pt]
  cz-slide & \ufig[0.48]{pv0-cz-slide} & semi-CX↑-CZ↓ & \ufig[0.48]{pv0-semi-CXup-CZdn}\\
\end{tabular}
\captionof{figure}{The sixteen rules of \protect\aid{{\_}CRel,{\_}==={\_}}
(\S\ref{sec:background}) as circuit equations, drawn from the Agda
terms.  Diagrams read in matrix-multiplication order, so a word $w
\bullet v$ is drawn with $v$ on the left.  \protect\aid{blake-c12}
relates two circuits implementing controlled-$Z$;
\protect\aid{yang-baxter} and \protect\aid{cz-slide} are the swap
calculus.  With the scalar restored (\S\ref{sec:composing}) only
\protect\aid{order-SH} changes, to $(SH)^3 = \omega^{-1/8}$.}
\Description{The sixteen circuit equations of the rule set.}
\label{fig:rules}
\end{center}

\begin{center}\footnotesize
  \begin{tabular}{@{}r@{\;\,}r@{\,}c@{\,}l@{\qquad}r@{\;\,}r@{\,}c@{\,}l@{}}
    $A_{a,b}$, $a\neq0$ & \ufig[0.58]{box-A-l} & $=$ & \ufig[0.58]{box-A-r} &
    $D_{a,b}$, $a\neq0$ & \ufig[0.58]{box-D-l} & $=$ & \ufig[0.58]{box-D-r}\\
    $A_{0,b}$, $b\neq0$ & \ufig[0.58]{box-A0-l} & $=$ & \ufig[0.58]{box-A0-r} &
    $D_{0,b}$ & \ufig[0.58]{box-D0-l} & $=$ & \ufig[0.58]{box-D0-r}\\
    $B_{a,b}$, $a\neq0$ & \ufig[0.58]{box-B-l} & $=$ & \ufig[0.58]{box-B-r} &
    $E_{b}$ & \ufig[0.58]{box-E-l} & $=$ & \ufig[0.58]{box-E-r}\\
    $B_{0,b}$ & \ufig[0.58]{box-B0-l} & $=$ & \ufig[0.58]{box-B0-r} & & & & \\
  \end{tabular}
\captionof{figure}{The four kinds of normal box (the qupit paper's Figure~6) and
how each is spelled as a circuit by the section of the normal form
(\S\ref{sec:symplectic}).}
\Description{The seven normal-box shapes and their circuit spellings.}
\label{fig:boxes}
\end{center}

\begin{center}
\resizebox{0.95\textwidth}{!}{%
\begin{tikzpicture}[x=1cm,y=1cm,
  box/.style={draw,rounded corners,align=center,font=\scriptsize,inner sep=2.5pt},
  lab/.style={font=\tiny,align=center,inner sep=1pt},
  >={Stealth[length=4pt]}]
  \node[box] (sp)    at (0,0)     {$\Sp$: full rules (17)\\coset tower, unique NF};
  \node[box] (sim)   at (3.35,0)  {$\Sp$: simplified\\rules (15)};
  \node[box] (sd)    at (6.7,0)   {$\PPauli\rtimes\Sp$: semidirect\\relation on $X,Z,H,S,CZ$};
  \node[box] (v1)    at (10.3,0)  {Simplified-V1\\(19 rules, $H,S,CZ$)};
  \node[box] (v0)    at (13.5,0)  {Paper-V0 (16 rules):\\our Figure~1 mod scalars};
  \node[box] (pauli) at (6.7,-1.5){$\PPauli$: $X^p,\ Z^p,\ XZ{=}ZX$};
  \node[box] (ex)    at (10.3,-1.5){Exact: Paper-V0 $+\ \omega$,\\one correction};
  \node[box] (pv1)   at (13.5,-1.5){Paper-V1 (15 rules):\\Figure~1 mod scalars};
  \draw[->] (sp) -- node[lab,above]{identity iso} (sim);
  \draw[->] (sim) -- node[lab,above]{semidirect thm} (sd);
  \draw[->] (pauli) -- node[lab,right]{semidirect thm} (sd);
  \draw[->] (sd) -- node[lab,above]{generator reduction:\\$X,Z \mapsto$ words} (v1);
  \draw[->] (v1) -- node[lab,above]{identity iso} (v0);
  \draw[->] (v0) -- node[lab,right]{identity iso} (pv1);
  \draw[->] (v0.south west) -- node[lab,below,pos=0.55]{central extension thm} (ex.north east);
\end{tikzpicture}}
\captionof{figure}{The chain of presentations: every arrow is a verified
presentation theorem or isomorphism of presentations; the semantic
group is $\Sp$ on the left, $\PPauli \rtimes \Sp$ from the semidirect
level on, and a central extension of it by $\langle\omega\rangle$ at
``Exact''.}
\Description{Boxes for the presentations, joined by arrows labelled with the theorem that connects them.}
\label{fig:route}
\end{center}

\begin{center}\begin{minipage}{\textwidth}\centering
\captionof{table}{Size of the development (lines of Agda incl.\ comments,
\texttt{wc -l}, September 2026).}
\label{tab:stats}
\small
\begin{tabular}{@{}lrr@{}}
\toprule
Component & Files & Lines \\
\midrule
Framework (\texttt{Word}, \texttt{Presentation}, \texttt{Circuit}, \texttt{Normalization}, \texttt{ForStdlib}) & 54 & 14\,848 \\
Symplectic group (normalisation 8\,352; box relations 18\,161; soundness etc.\ 8\,886) & 118 & 35\,399 \\
Projective Pauli group & 3 & 1\,220 \\
Projective Clifford: semidirect, plumbing, shared calculus & 26 & 17\,996 \\
Clifford with scalars & 11 & 5\,029 \\
\midrule
Other examples in the tree (symmetric, cyclic, wreath, Clifford+$T$, qubit Clifford) & 107 & 27\,552 \\
Not on the chain (\texttt{Simplified-V2}), index, notations & 10 & 4\,344 \\
\midrule
Total & 329 & 106\,388 \\
\bottomrule
\end{tabular}
\end{minipage}
\end{center}

\section{Selected Agda Definitions}\label{app:agda}

The normal-form datatype; the five hypotheses against which the library verifies a coset table
(the general form of the soundness law \aid{ract-sound} of
Appendix~\ref{sec:permutations}), in the generic \texttt{Transfer} module; the
six one-qupit rules of the simplified symplectic rule set
(\S\ref{sec:symplectic}), which with \aid{order-CZ}, \aid{comm-CZ-S↓},
\aid{comm-CZ-S↑} and \aid{selinger-c10}--\aid{c15} make its fifteen; and
the conjugation action of the semidirect presentation
(\S\ref{sec:composing}), the paper's Lemmas~2.22--2.23 as words.  All
three are quoted from the development, as is the doubly inductive
normal-form datatype of \S\ref{sec:symplectic}, shown first.

\begingroup\renewcommand{\AgdaCodeStyle}{\scriptsize}
\noindent\begin{minipage}[t]{0.44\textwidth}
\begin{code}%
\>[0]\AgdaFunction{M}\AgdaSpace{}%
\AgdaFunction{L}\AgdaSpace{}%
\AgdaFunction{ML}\AgdaSpace{}%
\AgdaSymbol{:}\AgdaSpace{}%
\AgdaDatatype{ℕ}\AgdaSpace{}%
\AgdaSymbol{→}\AgdaSpace{}%
\AgdaPrimitive{Set}\<%
\\
\>[0]\AgdaFunction{M}\AgdaSpace{}%
\AgdaNumber{0}\AgdaSpace{}%
\AgdaSymbol{=}\AgdaSpace{}%
\AgdaRecord{⊤}\<%
\\
\>[0]\AgdaFunction{M}\AgdaSpace{}%
\AgdaSymbol{(}\AgdaInductiveConstructor{₁₊}\AgdaSpace{}%
\AgdaBound{n}\AgdaSymbol{)}\AgdaSpace{}%
\AgdaSymbol{=}\AgdaSpace{}%
\AgdaDatatype{Vec}\AgdaSpace{}%
\AgdaPostulate{D}\AgdaSpace{}%
\AgdaBound{n}\AgdaSpace{}%
\AgdaOperator{\AgdaFunction{×}}\AgdaSpace{}%
\AgdaPostulate{E}\<%
\\
\>[0]\AgdaFunction{L}\AgdaSpace{}%
\AgdaNumber{0}\AgdaSpace{}%
\AgdaSymbol{=}\AgdaSpace{}%
\AgdaRecord{⊤}\<%
\\
\>[0]\AgdaFunction{L}\AgdaSpace{}%
\AgdaSymbol{(}\AgdaInductiveConstructor{₁₊}\AgdaSpace{}%
\AgdaBound{n}\AgdaSymbol{)}\AgdaSpace{}%
\AgdaSymbol{=}\AgdaSpace{}%
\AgdaDatatype{Vec}\AgdaSpace{}%
\AgdaPostulate{B}\AgdaSpace{}%
\AgdaBound{n}\AgdaSpace{}%
\AgdaOperator{\AgdaFunction{×}}\AgdaSpace{}%
\AgdaPostulate{A}\<%
\\
\>[0]\AgdaFunction{ML}\AgdaSpace{}%
\AgdaNumber{0}\AgdaSpace{}%
\AgdaSymbol{=}\AgdaSpace{}%
\AgdaRecord{⊤}\<%
\\
\>[0]\AgdaFunction{ML}\AgdaSpace{}%
\AgdaNumber{1}\AgdaSpace{}%
\AgdaSymbol{=}\AgdaSpace{}%
\AgdaFunction{M}\AgdaSpace{}%
\AgdaNumber{1}\AgdaSpace{}%
\AgdaOperator{\AgdaFunction{×}}\AgdaSpace{}%
\AgdaFunction{L}\AgdaSpace{}%
\AgdaNumber{1}\<%
\\
\>[0]\AgdaFunction{ML}\AgdaSpace{}%
\AgdaBound{m}\AgdaSymbol{@(}\AgdaInductiveConstructor{₂₊}\AgdaSpace{}%
\AgdaBound{n}\AgdaSymbol{)}%
\>[350I]\AgdaSymbol{=}\AgdaSpace{}%
\AgdaFunction{M}\AgdaSpace{}%
\AgdaBound{m}\AgdaSpace{}%
\AgdaOperator{\AgdaFunction{×}}\AgdaSpace{}%
\AgdaFunction{L}\AgdaSpace{}%
\AgdaBound{m}\<%
\\
\>[.][@{}l@{}]\<[350I]%
\>[12]\AgdaOperator{\AgdaDatatype{⊎}}\AgdaSpace{}%
\AgdaPostulate{D}\AgdaSpace{}%
\AgdaOperator{\AgdaFunction{×}}\AgdaSpace{}%
\AgdaFunction{ML}\AgdaSpace{}%
\AgdaSymbol{(}\AgdaInductiveConstructor{₁₊}\AgdaSpace{}%
\AgdaBound{n}\AgdaSymbol{)}\<%
\end{code}

\end{minipage}\hfill
\begin{minipage}[t]{0.54\textwidth}
\begin{code}%
\>[0]\AgdaFunction{NF}\AgdaSpace{}%
\AgdaSymbol{:}\AgdaSpace{}%
\AgdaDatatype{ℕ}\AgdaSpace{}%
\AgdaSymbol{→}\AgdaSpace{}%
\AgdaPrimitive{Set}\<%
\\
\>[0]\AgdaFunction{NF}\AgdaSpace{}%
\AgdaNumber{0}\AgdaSpace{}%
\AgdaSymbol{=}\AgdaSpace{}%
\AgdaRecord{⊤}\<%
\\
\>[0]\AgdaFunction{NF}\AgdaSpace{}%
\AgdaSymbol{(}\AgdaInductiveConstructor{₁₊}\AgdaSpace{}%
\AgdaBound{n}\AgdaSymbol{)}\AgdaSpace{}%
\AgdaSymbol{=}\AgdaSpace{}%
\AgdaFunction{NF}\AgdaSpace{}%
\AgdaBound{n}\AgdaSpace{}%
\AgdaOperator{\AgdaFunction{×}}\AgdaSpace{}%
\AgdaFunction{ML}\AgdaSpace{}%
\AgdaSymbol{(}\AgdaInductiveConstructor{₁₊}\AgdaSpace{}%
\AgdaBound{n}\AgdaSymbol{)}\<%
\\
%
\\[\AgdaEmptyExtraSkip]%
\>[0]\AgdaOperator{\AgdaFunction{[\AgdaUnderscore{}]}}\AgdaSpace{}%
\AgdaSymbol{:}\AgdaSpace{}%
\AgdaSymbol{∀}\AgdaSpace{}%
\AgdaSymbol{\{}\AgdaBound{n}\AgdaSymbol{\}}\AgdaSpace{}%
\AgdaSymbol{→}\AgdaSpace{}%
\AgdaFunction{NF}\AgdaSpace{}%
\AgdaBound{n}\AgdaSpace{}%
\AgdaSymbol{→}\AgdaSpace{}%
\AgdaDatatype{Word}\AgdaSpace{}%
\AgdaSymbol{(}\AgdaDatatype{Gen}\AgdaSpace{}%
\AgdaBound{n}\AgdaSymbol{)}\<%
\\
\>[0]\AgdaOperator{\AgdaFunction{[\AgdaUnderscore{}]}}\AgdaSpace{}%
\AgdaSymbol{\{}\AgdaNumber{0}\AgdaSymbol{\}}%
\>[12]\AgdaInductiveConstructor{tt}%
\>[23]\AgdaSymbol{=}\AgdaSpace{}%
\AgdaInductiveConstructor{ε}\<%
\\
\>[0]\AgdaOperator{\AgdaFunction{[\AgdaUnderscore{}]}}\AgdaSpace{}%
\AgdaSymbol{\{}\AgdaInductiveConstructor{₁₊}\AgdaSpace{}%
\AgdaBound{n}\AgdaSymbol{\}}%
\>[12]\AgdaSymbol{(}\AgdaBound{nf}\AgdaSpace{}%
\AgdaOperator{\AgdaInductiveConstructor{,}}\AgdaSpace{}%
\AgdaBound{ml}\AgdaSymbol{)}%
\>[23]\AgdaSymbol{=}\AgdaSpace{}%
\AgdaOperator{\AgdaFunction{[}}\AgdaSpace{}%
\AgdaBound{nf}\AgdaSpace{}%
\AgdaOperator{\AgdaFunction{]}}\AgdaSpace{}%
\AgdaOperator{\AgdaPostulate{↑}}\AgdaSpace{}%
\AgdaOperator{\AgdaInductiveConstructor{•}}\AgdaSpace{}%
\AgdaOperator{\AgdaPostulate{[}}\AgdaSpace{}%
\AgdaBound{ml}\AgdaSpace{}%
\AgdaOperator{\AgdaPostulate{]ᵐˡ}}\<%
\end{code}

\end{minipage}
%% Copied from ../vqupit/latex/vqupit/sections/framework.tex (framework, block 2), the LaTeX
%% that `agda --latex --only-scope-checking' generated for the long version.
\begin{code}%
\>[0][@{}l@{\AgdaIndent{1}}]%
\>[2]\AgdaSymbol{(}\AgdaBound{h=⁻¹f-gen}\AgdaSpace{}%
\AgdaSymbol{:}\AgdaSpace{}%
\AgdaSymbol{∀}\AgdaSpace{}%
\AgdaSymbol{(}\AgdaBound{x}\AgdaSpace{}%
\AgdaSymbol{:}\AgdaSpace{}%
\AgdaPostulate{X}\AgdaSymbol{)}\AgdaSpace{}%
\AgdaSymbol{→}\AgdaSpace{}%
\AgdaSymbol{(}\AgdaOperator{\AgdaInductiveConstructor{[}}\AgdaSpace{}%
\AgdaBound{x}\AgdaSpace{}%
\AgdaOperator{\AgdaInductiveConstructor{]ʷ}}\AgdaSpace{}%
\AgdaOperator{\AgdaInductiveConstructor{,}}\AgdaSpace{}%
\AgdaPostulate{I}\AgdaSymbol{)}\AgdaSpace{}%
\AgdaOperator{\AgdaPostulate{\textasciitilde{}}}\AgdaSpace{}%
\AgdaSymbol{((}\AgdaPostulate{h}\AgdaSpace{}%
\AgdaOperator{\AgdaFunction{ᵗ}}\AgdaSymbol{)}\AgdaSpace{}%
\AgdaPostulate{I}\AgdaSpace{}%
\AgdaSymbol{(}\AgdaPostulate{f}\AgdaSpace{}%
\AgdaBound{x}\AgdaSymbol{)))}%
\>[84]\AgdaComment{--\ 1}\<%
\\
%
\>[2]\AgdaSymbol{(}\AgdaBound{h-wd-ax}\AgdaSpace{}%
\AgdaSymbol{:}\AgdaSpace{}%
\AgdaSymbol{∀}\AgdaSpace{}%
\AgdaSymbol{(}\AgdaBound{c}\AgdaSpace{}%
\AgdaSymbol{:}\AgdaSpace{}%
\AgdaPostulate{C}\AgdaSymbol{)\{}\AgdaBound{u}\AgdaSpace{}%
\AgdaBound{t}\AgdaSpace{}%
\AgdaSymbol{:}\AgdaSpace{}%
\AgdaDatatype{Word}\AgdaSpace{}%
\AgdaPostulate{Y}\AgdaSymbol{\}}\AgdaSpace{}%
\AgdaSymbol{→}\AgdaSpace{}%
\AgdaBound{u}\AgdaSpace{}%
\AgdaOperator{\AgdaPostulate{===₂}}\AgdaSpace{}%
\AgdaBound{t}\AgdaSpace{}%
\AgdaSymbol{→}\AgdaSpace{}%
\AgdaSymbol{((}\AgdaPostulate{h}\AgdaSpace{}%
\AgdaOperator{\AgdaFunction{ᵗ}}\AgdaSymbol{)}\AgdaSpace{}%
\AgdaBound{c}\AgdaSpace{}%
\AgdaBound{u}\AgdaSymbol{)}\AgdaSpace{}%
\AgdaOperator{\AgdaPostulate{\textasciitilde{}}}\AgdaSpace{}%
\AgdaSymbol{((}\AgdaPostulate{h}\AgdaSpace{}%
\AgdaOperator{\AgdaFunction{ᵗ}}\AgdaSymbol{)}\AgdaSpace{}%
\AgdaBound{c}\AgdaSpace{}%
\AgdaBound{t}\AgdaSymbol{))}%
\>[84]\AgdaComment{--\ 2}\<%
\\
%
\>[2]\AgdaSymbol{(}\AgdaBound{f-wd-ax}\AgdaSpace{}%
\AgdaSymbol{:}\AgdaSpace{}%
\AgdaSymbol{∀}\AgdaSpace{}%
\AgdaSymbol{\{}\AgdaBound{w}\AgdaSpace{}%
\AgdaBound{v}\AgdaSymbol{\}}\AgdaSpace{}%
\AgdaSymbol{→}\AgdaSpace{}%
\AgdaBound{w}\AgdaSpace{}%
\AgdaOperator{\AgdaPostulate{===₁}}\AgdaSpace{}%
\AgdaBound{v}\AgdaSpace{}%
\AgdaSymbol{→}\AgdaSpace{}%
\AgdaSymbol{(}\AgdaPostulate{f}\AgdaSpace{}%
\AgdaOperator{\AgdaFunction{ʷ}}\AgdaSymbol{)}\AgdaSpace{}%
\AgdaBound{w}\AgdaSpace{}%
\AgdaOperator{\AgdaPostulate{≈₂}}\AgdaSpace{}%
\AgdaSymbol{(}\AgdaPostulate{f}\AgdaSpace{}%
\AgdaOperator{\AgdaFunction{ʷ}}\AgdaSymbol{)}\AgdaSpace{}%
\AgdaBound{v}\AgdaSymbol{)}%
\>[84]\AgdaComment{--\ 3}\<%
\\
%
\>[2]\AgdaSymbol{(}\AgdaBound{[I]≈ε}\AgdaSpace{}%
\AgdaSymbol{:}\AgdaSpace{}%
\AgdaOperator{\AgdaPostulate{[}}\AgdaSpace{}%
\AgdaPostulate{I}\AgdaSpace{}%
\AgdaOperator{\AgdaPostulate{]}}\AgdaSpace{}%
\AgdaOperator{\AgdaPostulate{≈₂}}\AgdaSpace{}%
\AgdaInductiveConstructor{ε}\AgdaSymbol{)}%
\>[84]\AgdaComment{--\ 4}\<%
\\
%
\>[2]\AgdaSymbol{(}\AgdaBound{h=ract}\AgdaSpace{}%
\AgdaSymbol{:}\AgdaSpace{}%
\AgdaSymbol{∀}\AgdaSpace{}%
\AgdaBound{c}\AgdaSpace{}%
\AgdaBound{b}\AgdaSpace{}%
\AgdaSymbol{→}\AgdaSpace{}%
\AgdaKeyword{let}\AgdaSpace{}%
\AgdaSymbol{(}\AgdaBound{b'}\AgdaSpace{}%
\AgdaOperator{\AgdaInductiveConstructor{,}}\AgdaSpace{}%
\AgdaBound{c'}\AgdaSymbol{)}\AgdaSpace{}%
\AgdaSymbol{=}\AgdaSpace{}%
\AgdaPostulate{h}\AgdaSpace{}%
\AgdaBound{c}\AgdaSpace{}%
\AgdaBound{b}\AgdaSpace{}%
\AgdaKeyword{in}\AgdaSpace{}%
\AgdaOperator{\AgdaPostulate{[}}\AgdaSpace{}%
\AgdaBound{c}\AgdaSpace{}%
\AgdaOperator{\AgdaPostulate{]}}\AgdaSpace{}%
\AgdaOperator{\AgdaInductiveConstructor{•}}\AgdaSpace{}%
\AgdaOperator{\AgdaInductiveConstructor{[}}\AgdaSpace{}%
\AgdaBound{b}\AgdaSpace{}%
\AgdaOperator{\AgdaInductiveConstructor{]ʷ}}\AgdaSpace{}%
\AgdaOperator{\AgdaPostulate{≈₂}}\AgdaSpace{}%
\AgdaSymbol{(}\AgdaPostulate{f}\AgdaSpace{}%
\AgdaOperator{\AgdaFunction{ʷ}}\AgdaSymbol{)}\AgdaSpace{}%
\AgdaBound{b'}\AgdaSpace{}%
\AgdaOperator{\AgdaInductiveConstructor{•}}\AgdaSpace{}%
\AgdaOperator{\AgdaPostulate{[}}\AgdaSpace{}%
\AgdaBound{c'}\AgdaSpace{}%
\AgdaOperator{\AgdaPostulate{]}}\AgdaSymbol{)}%
\>[84]\AgdaComment{--\ 5}\<%
\end{code}

\noindent\begin{minipage}[t]{0.5\textwidth}
\begin{code}%
\>[0]\AgdaFunction{order-S}\AgdaSpace{}%
\AgdaSymbol{:}%
\>[20]\AgdaSymbol{(}\AgdaInductiveConstructor{₁₊}\AgdaSpace{}%
\AgdaGeneralizable{n}\AgdaSymbol{)}\AgdaSpace{}%
\AgdaOperator{\AgdaPostulate{SRel,}}%
\>[34]\AgdaPostulate{S}\AgdaSpace{}%
\AgdaOperator{\AgdaFunction{\textasciicircum{}}}\AgdaSpace{}%
\AgdaPostulate{p}\AgdaSpace{}%
\AgdaOperator{\AgdaPostulate{===}}\AgdaSpace{}%
\AgdaInductiveConstructor{ε}\<%
\\
\>[0]\AgdaFunction{order-H}\AgdaSpace{}%
\AgdaSymbol{:}%
\>[20]\AgdaSymbol{(}\AgdaInductiveConstructor{₁₊}\AgdaSpace{}%
\AgdaGeneralizable{n}\AgdaSymbol{)}\AgdaSpace{}%
\AgdaOperator{\AgdaPostulate{SRel,}}%
\>[34]\AgdaPostulate{H}\AgdaSpace{}%
\AgdaOperator{\AgdaFunction{\textasciicircum{}}}\AgdaSpace{}%
\AgdaNumber{2}\AgdaSpace{}%
\AgdaOperator{\AgdaPostulate{===}}\AgdaSpace{}%
\AgdaPostulate{M₋₁}\<%
\\
\>[0]\AgdaFunction{M-power}\AgdaSpace{}%
\AgdaSymbol{:}%
\>[14]\AgdaSymbol{∀}\AgdaSpace{}%
\AgdaBound{k}\AgdaSpace{}%
\AgdaSymbol{→}\AgdaSpace{}%
\AgdaSymbol{(}\AgdaInductiveConstructor{₁₊}\AgdaSpace{}%
\AgdaGeneralizable{n}\AgdaSymbol{)}\AgdaSpace{}%
\AgdaOperator{\AgdaPostulate{SRel,}}%
\>[34]\AgdaPostulate{Mg\textasciicircum{}}\AgdaSpace{}%
\AgdaBound{k}\AgdaSpace{}%
\AgdaOperator{\AgdaPostulate{===}}\AgdaSpace{}%
\AgdaFunction{M}\AgdaSpace{}%
\AgdaSymbol{(}\AgdaPostulate{g\textasciicircum{}}\AgdaSpace{}%
\AgdaBound{k}\AgdaSymbol{)}\<%
\\
\>[0]\AgdaFunction{semi-MS}\AgdaSpace{}%
\AgdaSymbol{:}%
\>[20]\AgdaSymbol{(}\AgdaInductiveConstructor{₁₊}\AgdaSpace{}%
\AgdaGeneralizable{n}\AgdaSymbol{)}\AgdaSpace{}%
\AgdaOperator{\AgdaPostulate{SRel,}}%
\>[34]\AgdaPostulate{Mg}\AgdaSpace{}%
\AgdaOperator{\AgdaInductiveConstructor{•}}\AgdaSpace{}%
\AgdaPostulate{S}\AgdaSpace{}%
\AgdaOperator{\AgdaPostulate{===}}\AgdaSpace{}%
\AgdaPostulate{S\textasciicircum{}}\AgdaSpace{}%
\AgdaSymbol{(}\AgdaPostulate{g}\AgdaSpace{}%
\AgdaOperator{\AgdaPostulate{*}}\AgdaSpace{}%
\AgdaPostulate{g}\AgdaSymbol{)}\AgdaSpace{}%
\AgdaOperator{\AgdaInductiveConstructor{•}}\AgdaSpace{}%
\AgdaPostulate{Mg}\<%
\\
\>[0]\AgdaFunction{semi-M↑CZ}\AgdaSpace{}%
\AgdaSymbol{:}%
\>[20]\AgdaSymbol{(}\AgdaInductiveConstructor{₂₊}\AgdaSpace{}%
\AgdaGeneralizable{n}\AgdaSymbol{)}\AgdaSpace{}%
\AgdaOperator{\AgdaPostulate{SRel,}}%
\>[34]\AgdaPostulate{Mg}\AgdaSpace{}%
\AgdaOperator{\AgdaPostulate{↑}}\AgdaSpace{}%
\AgdaOperator{\AgdaInductiveConstructor{•}}\AgdaSpace{}%
\AgdaPostulate{CZ}\AgdaSpace{}%
\AgdaOperator{\AgdaPostulate{===}}\AgdaSpace{}%
\AgdaPostulate{CZ\textasciicircum{}}\AgdaSpace{}%
\AgdaPostulate{g}\AgdaSpace{}%
\AgdaOperator{\AgdaInductiveConstructor{•}}\AgdaSpace{}%
\AgdaPostulate{Mg}\AgdaSpace{}%
\AgdaOperator{\AgdaPostulate{↑}}\<%
\\
\>[0]\AgdaFunction{semi-M↓CZ}\AgdaSpace{}%
\AgdaSymbol{:}%
\>[20]\AgdaSymbol{(}\AgdaInductiveConstructor{₂₊}\AgdaSpace{}%
\AgdaGeneralizable{n}\AgdaSymbol{)}\AgdaSpace{}%
\AgdaOperator{\AgdaPostulate{SRel,}}%
\>[34]\AgdaPostulate{Mg}\AgdaSpace{}%
\AgdaOperator{\AgdaPostulate{↓}}\AgdaSpace{}%
\AgdaOperator{\AgdaInductiveConstructor{•}}\AgdaSpace{}%
\AgdaPostulate{CZ}\AgdaSpace{}%
\AgdaOperator{\AgdaPostulate{===}}\AgdaSpace{}%
\AgdaPostulate{CZ\textasciicircum{}}\AgdaSpace{}%
\AgdaPostulate{g}\AgdaSpace{}%
\AgdaOperator{\AgdaInductiveConstructor{•}}\AgdaSpace{}%
\AgdaPostulate{Mg}\AgdaSpace{}%
\AgdaOperator{\AgdaPostulate{↓}}\<%
\end{code}

\end{minipage}\hfill
\begin{minipage}[t]{0.47\textwidth}
\begin{code}%
\>[0]\AgdaFunction{conj}\AgdaSpace{}%
\AgdaSymbol{:}\AgdaSpace{}%
\AgdaDatatype{Sym.Gen}\AgdaSpace{}%
\AgdaGeneralizable{n}\AgdaSpace{}%
\AgdaSymbol{→}\AgdaSpace{}%
\AgdaDatatype{XZ.Gen}\AgdaSpace{}%
\AgdaGeneralizable{n}\AgdaSpace{}%
\AgdaSymbol{→}\AgdaSpace{}%
\AgdaDatatype{Word}\AgdaSpace{}%
\AgdaSymbol{(}\AgdaDatatype{XZ.Gen}\AgdaSpace{}%
\AgdaGeneralizable{n}\AgdaSymbol{)}\<%
\\
\>[0]\AgdaFunction{conj}\AgdaSpace{}%
\AgdaInductiveConstructor{Sym.H-gen}%
\>[17]\AgdaInductiveConstructor{X-gen}%
\>[33]\AgdaSymbol{=}\AgdaSpace{}%
\AgdaPostulate{Z}\<%
\\
\>[0]\AgdaFunction{conj}\AgdaSpace{}%
\AgdaInductiveConstructor{Sym.H-gen}%
\>[17]\AgdaInductiveConstructor{Z-gen}%
\>[33]\AgdaSymbol{=}\AgdaSpace{}%
\AgdaPostulate{X}\AgdaSpace{}%
\AgdaOperator{\AgdaFunction{\textasciicircum{}}}\AgdaSpace{}%
\AgdaPostulate{p-1}\<%
\\
\>[0]\AgdaFunction{conj}\AgdaSpace{}%
\AgdaInductiveConstructor{Sym.S-gen}%
\>[17]\AgdaInductiveConstructor{X-gen}%
\>[33]\AgdaSymbol{=}\AgdaSpace{}%
\AgdaPostulate{X}\AgdaSpace{}%
\AgdaOperator{\AgdaInductiveConstructor{•}}\AgdaSpace{}%
\AgdaPostulate{Z}\<%
\\
\>[0]\AgdaFunction{conj}\AgdaSpace{}%
\AgdaInductiveConstructor{Sym.S-gen}%
\>[17]\AgdaInductiveConstructor{Z-gen}%
\>[33]\AgdaSymbol{=}\AgdaSpace{}%
\AgdaPostulate{Z}\<%
\\
\>[0]\AgdaFunction{conj}\AgdaSpace{}%
\AgdaInductiveConstructor{Sym.CZ-gen}%
\>[17]\AgdaInductiveConstructor{X-gen}%
\>[33]\AgdaSymbol{=}\AgdaSpace{}%
\AgdaPostulate{X}\AgdaSpace{}%
\AgdaOperator{\AgdaInductiveConstructor{•}}\AgdaSpace{}%
\AgdaPostulate{Z}\AgdaSpace{}%
\AgdaOperator{\AgdaPostulate{↑}}\<%
\\
\>[0]\AgdaFunction{conj}\AgdaSpace{}%
\AgdaInductiveConstructor{Sym.CZ-gen}%
\>[17]\AgdaInductiveConstructor{Z-gen}%
\>[33]\AgdaSymbol{=}\AgdaSpace{}%
\AgdaPostulate{Z}\<%
\\
\>[0]\AgdaFunction{conj}\AgdaSpace{}%
\AgdaInductiveConstructor{Sym.CZ-gen}%
\>[17]\AgdaSymbol{(}\AgdaInductiveConstructor{X-gen}\AgdaSpace{}%
\AgdaOperator{\AgdaInductiveConstructor{↥}}\AgdaSymbol{)}%
\>[33]\AgdaSymbol{=}\AgdaSpace{}%
\AgdaPostulate{X}\AgdaSpace{}%
\AgdaOperator{\AgdaPostulate{↑}}\AgdaSpace{}%
\AgdaOperator{\AgdaInductiveConstructor{•}}\AgdaSpace{}%
\AgdaPostulate{Z}\<%
\\
\>[0]\AgdaFunction{conj}\AgdaSpace{}%
\AgdaInductiveConstructor{Sym.CZ-gen}%
\>[17]\AgdaSymbol{(}\AgdaInductiveConstructor{Z-gen}\AgdaSpace{}%
\AgdaOperator{\AgdaInductiveConstructor{↥}}\AgdaSymbol{)}%
\>[33]\AgdaSymbol{=}\AgdaSpace{}%
\AgdaPostulate{Z}\AgdaSpace{}%
\AgdaOperator{\AgdaPostulate{↑}}\<%
\\
\>[0]\AgdaFunction{conj}\AgdaSpace{}%
\AgdaInductiveConstructor{Sym.H-gen}%
\>[17]\AgdaSymbol{(}\AgdaBound{xz}\AgdaSpace{}%
\AgdaOperator{\AgdaInductiveConstructor{↥}}\AgdaSymbol{)}%
\>[33]\AgdaSymbol{=}\AgdaSpace{}%
\AgdaOperator{\AgdaInductiveConstructor{[}}\AgdaSpace{}%
\AgdaBound{xz}\AgdaSpace{}%
\AgdaOperator{\AgdaInductiveConstructor{]ʷ}}\AgdaSpace{}%
\AgdaOperator{\AgdaPostulate{↑}}\<%
\\
\>[0]\AgdaFunction{conj}\AgdaSpace{}%
\AgdaInductiveConstructor{Sym.S-gen}%
\>[17]\AgdaSymbol{(}\AgdaBound{xz}\AgdaSpace{}%
\AgdaOperator{\AgdaInductiveConstructor{↥}}\AgdaSymbol{)}%
\>[33]\AgdaSymbol{=}\AgdaSpace{}%
\AgdaOperator{\AgdaInductiveConstructor{[}}\AgdaSpace{}%
\AgdaBound{xz}\AgdaSpace{}%
\AgdaOperator{\AgdaInductiveConstructor{]ʷ}}\AgdaSpace{}%
\AgdaOperator{\AgdaPostulate{↑}}\<%
\\
\>[0]\AgdaFunction{conj}\AgdaSpace{}%
\AgdaInductiveConstructor{Sym.CZ-gen}%
\>[17]\AgdaSymbol{(}\AgdaBound{xz}\AgdaSpace{}%
\AgdaOperator{\AgdaInductiveConstructor{↥}}\AgdaSpace{}%
\AgdaOperator{\AgdaInductiveConstructor{↥}}\AgdaSymbol{)}%
\>[33]\AgdaSymbol{=}\AgdaSpace{}%
\AgdaOperator{\AgdaInductiveConstructor{[}}\AgdaSpace{}%
\AgdaBound{xz}\AgdaSpace{}%
\AgdaOperator{\AgdaInductiveConstructor{]ʷ}}\AgdaSpace{}%
\AgdaOperator{\AgdaPostulate{↑}}\AgdaSpace{}%
\AgdaOperator{\AgdaPostulate{↑}}\<%
\\
\>[0]\AgdaFunction{conj}\AgdaSpace{}%
\AgdaSymbol{(}\AgdaBound{sym}\AgdaSpace{}%
\AgdaOperator{\AgdaInductiveConstructor{↥}}\AgdaSymbol{)}%
\>[17]\AgdaInductiveConstructor{X-gen}%
\>[33]\AgdaSymbol{=}\AgdaSpace{}%
\AgdaPostulate{X}\<%
\\
\>[0]\AgdaFunction{conj}\AgdaSpace{}%
\AgdaSymbol{(}\AgdaBound{sym}\AgdaSpace{}%
\AgdaOperator{\AgdaInductiveConstructor{↥}}\AgdaSymbol{)}%
\>[17]\AgdaInductiveConstructor{Z-gen}%
\>[33]\AgdaSymbol{=}\AgdaSpace{}%
\AgdaPostulate{Z}\<%
\\
\>[0]\AgdaFunction{conj}\AgdaSpace{}%
\AgdaSymbol{(}\AgdaBound{sym}\AgdaSpace{}%
\AgdaOperator{\AgdaInductiveConstructor{↥}}\AgdaSymbol{)}%
\>[17]\AgdaSymbol{(}\AgdaBound{xz}\AgdaSpace{}%
\AgdaOperator{\AgdaInductiveConstructor{↥}}\AgdaSymbol{)}%
\>[33]\AgdaSymbol{=}\AgdaSpace{}%
\AgdaFunction{conj}\AgdaSpace{}%
\AgdaBound{sym}\AgdaSpace{}%
\AgdaBound{xz}\AgdaSpace{}%
\AgdaOperator{\AgdaPostulate{↑}}\<%
\end{code}

\end{minipage}
\endgroup
